
\ifintegriereaufgaben

\section*{Reihen Aufgaben}\label{MA2-zahlenreihen-aufgaben}
Erkl�rungen und Beispiele, die zur Bearbeitung der Aufgaben hilfreich sind, finden Sie im Skript\wennreferenzvorhanden{MA2-zahlenreihen}{ (Abschnitt \abschnittsverweis{MA2-zahlenreihen})} und, falls Sie nicht weiterkommen, dann schauen Sie \hyperref[MA2-zahlenreihen-loesungen]{hier}.
\else
\section{Reihen Aufgaben}\label{MA2-zahlenreihen-aufgaben}

Regeln und Beispiele zu diesen Aufgaben finden Sie im Skript. Die L�sungen finden Sie \hyperref[MA2-zahlenreihen-loesungen]{hier}.
\fi

\begin{exercisebox}[Reihen] %A.Z. LS
	Geben sind die folgenden Reihen
	\begin{eqnarray*}
		&(a)\,& 1+1+1+1 + \dots \\
		&(b)\,& 1-1+1-1 + \dots
	\end{eqnarray*}
	Schreiben Sie die Reihen in der Form \( \displaystyle{\sum\limits_{k=1}^\infty} a_k\) und bestimmen Sie die Partialsummen $S_n$ f�r $n\in\mathbb N$ und gegebenenfalls den Wert der Reihe.
\end{exercisebox}


\begin{exercisebox}[Reihen] %Papula I, VI, Abschnitt 1, Aufgabe 1
	Berechnen Sie den Wert der folgenden Reihen
	\[
		(a)\, \sum\limits_{k=1}^\infty \left( \frac18\right)^{k-1}\quad\quad (b)\, \sum\limits_{k=1}^\infty 0,3^{k-1} \quad\quad (c)\,\sum\limits_{k=1}^\infty 4\cdot \left( - \frac23\right)^{k-1}
	\]
\end{exercisebox}

\begin{exercisebox}[Reihen] %HMI I, Blatt 9, Aufgabe 1
	Untersuchen Sie das Konvergenzverhalten der folgenden Reihen
	\[
		(a)\, \sum\limits_{k=0}^\infty \frac{k!}{k^k}\quad\quad\quad (b)\, \sum\limits_{k=0}^\infty \frac{k^4}{3^k} \quad\quad\quad (c)\,\sum\limits_{k=1}^\infty \frac{(k+1)^{k-1}}{(-k)^k}
	\]
\end{exercisebox}

\begin{exercisebox}[Rechenregeln] %HMI I, Blatt 9, Aufgabe 3
	Berechnen Sie die Werte der folgenden Reihen
	\[
		(a)\, \sum\limits_{k=0}^\infty \left( -\frac{9}{19}\right)^k\quad (b)\, \sum\limits_{k=2}^\infty \frac{3^k}{4^{k+1}} \quad (c)\,\sum\limits_{k=0}^\infty \frac{\sqrt{121} + 5^{2k-1}}{29^{k+1}}
	\]
\end{exercisebox}

\begin{exercisebox}[Konvergenz] %HMI I, Blatt 9, Aufgabe 2
	\begin{itemize}
		\item[(a)] F�r welche $x\in\mathbb R$ konvergiert die unendliche Reihe
			\[
				S(x):= \sum\limits_{k=1}^\infty \frac{x^{2k}}{(1-4|x|)^{k-1}}\textnormal{?} \]
				Hinweis: Klammern Sie eine geeignete Potenz von $x^k$ aus und l�sen Sie die Ungleichung mit GeoGebra o.�..
			\item[(b)] Bestimmen Sie f�r diejenigen $x\in\mathbb R$, f�r welche die Reihe $S(x)$ konvergiert, den Reihenwert.
		\end{itemize}
\end{exercisebox}


%\begin{exercisebox}[Konvergenz] %HMI I, Blatt 8, Aufgabe 2
%	Es sei \( (a_n)_{n\in\mathbb N}\) eine Folge, die gegen $a\in\mathbb R$ koknvergiert und \( (b_n)_{n\in\mathbb N}\) die Folge, deren $n$-tes Folgenglied
%	\[ b_n = \frac{1}{n+1}\sum\limits_{k=0}^n a_k \]
%	lautet.
%
%	Zeigen Sie, dass $(b_n)_{n\in\mathbb N}$ auch gegen $a$ konvergiert. Hinweis: Betrachten Sie zun�chst den Fall $a=0$.
%\vspace{0.3cm}
%{\bf L�sung:} \newline
%\end{exercisebox}

%\subsection{Praxis}

%\begin{mybox}[MC - Reihen]
%\begin{itemize}
%	\item Typen von Folgen (monoton, Nullfolge)
%	\item Typen von Reihen (alternierend, divergent)
%	\item notwenige und hinreichende Kriterien
%\end{itemize}
%\end{mybox}


